% 本文件由 LaTeX 排版 Agent 根据有效 translation.json 自动生成。
% author=Mathetes; work_id=0101; title=马特特斯致狄奥格内托书

\chapter{\WorkTitle{\Person{马特特斯}致\Person{狄奥格内托书}}}

% structure: p0001 source=h1 target=omitted reason=page-title-already-rendered-by-parent

% structure: p0002 source=h2 target=section reason=relative-heading-rank; base=h2

\section{书函的缘起}

鉴于我看到你，最卓越的\Person{狄奥格内托}，极其渴望了解基督徒中盛行的敬拜天主的方式，并非常仔细而恳切地询问他们，所信靠的是哪位天主，遵行何种宗教形式，以至于全都轻视世界本身，并轻看死亡，既不认为希腊人所尊崇的那些是神，也不持守犹太人的迷信；以及他们彼此之间怀有何种情感；最后，为何这种新的虔敬样式或做法直到如今才进入世界，而非很久以前；我衷心欢迎你的这个渴望，我恳求那使我们能言又能听的天主，赐予我如此发言，使我首先能听到你得到造就，并赐予你如此聆听，使我发言之人不致为此后悔。\par

% structure: p0004 source=h2 target=section reason=relative-heading-rank; base=h2

\section{偶像的虚妄}

那么，来吧，在摆脱了你心中所有偏见之后，放下你所习惯的、易于欺骗你的事物，并如同从起初成为新人，因为按你自己的承认，你将要聆听一种新教义的听众；来默观，不只以眼目，更以理智，你所宣称并视为神的那种实体的本质与形状。\par

它们的其中一个不就是一块与我们所踩踏的相似的石头吗？第二个不就是在任何方面都不优于为我们日常使用所制造的器皿的黄铜吗？第三个不就是木头，且已经朽坏吗？第四个不就是银子，需要人看守，免得被偷吗？第五个不就是铁，被锈所腐蚀吗？第六个不就是陶器，在任何程度上都不比为最卑微目的所造的更有价值吗？\par

所有这些不都是可朽坏的物质吗？它们不都是藉着铁和火制造的吗？雕刻匠不是造了其中一个，铜匠造了第二个，银匠造了第三个，陶匠造了第四个吗？在它们被这些工匠的技艺塑造成这些神明的形状之前，它们中的每一个不都是以各自的方式易于变化吗？那些现在由相同材料制成的器皿，若遇到相同的工匠，岂不也会变得像这些吗？这些现在被你们敬拜的，难道不能被人类重新制成与其它相似的器皿吗？它们不都是聋的吗？不都是瞎的吗？不都是没有生命的吗？不都是没有感觉的吗？不都是不能动的吗？不都是容易朽坏的吗？不都是可朽的吗？\par

这些你们称为神；这些你们事奉；这些你们敬拜；而你们完全变得像它们。为此你们恨恶基督徒，因为他们不认为\Emph{这些}是神。但难道你们自己，现在认为并假设这些是神的，不是比他们更加轻视它们吗？你们岂不是更多地嘲笑和侮辱它们，当你们敬拜那些由石头和陶器制成的，却不派人看守它们；而那些由金银制成的，你们夜间锁起来，白天派看守照管，免得被偷？而且藉着你们打算献给它们的那些礼物，如果它们有感觉，你们岂不是在惩罚而不是尊崇它们吗？但如果相反，它们没有感觉，你们就在以血和祭物的烟敬拜它们时，证明了这一事实。让你们中的任何一个人忍受这样的侮辱！让你们中的任何一个人容忍这样的事情加在自己身上！但没有一个人，除非被迫，会忍受这样的对待，因为他有感觉和理性。然而石头却甘心忍受，因为它没有感觉。当然，你们并未通过行为显示你们的神拥有感觉。至于基督徒不习惯事奉这样的神这一事实，我很容易找到许多其他可说的；但如果已经说的对任何人来说似乎都不足够，我认为再说任何话都是徒劳的。\par

% structure: p0009 source=h2 target=section reason=relative-heading-rank; base=h2

\section{犹太人的迷信}

其次，我想你非常渴望听到有关这一点的话，即基督徒不遵行与犹太人相同形式的敬拜。那么，犹太人，如果他们避免了上述的事奉方式，并认为敬拜独一的天主为万有的主是合适的，就是对的；但如果他们以我们所描述的方式敬拜祂，就大错了。因为外邦人将这些东西献给没有感觉和听觉的，是提供了疯狂的例证；而他们，却以为将这些东西献给天主，好像祂需要它们，正当地可以被视为愚蠢的行为，而非敬拜。因为那造了天和地及其中的万物，并将我们所需的一切赐给我们的，当然不需要那些祂自己所赐给那些想要献上给祂之人的任何东西。但那些想象藉着血、祭物的烟和燔祭，他们献上蒙悦纳的祭物给祂，并以为藉着这样的尊荣，他们表示对祂的尊敬——这些人，因设想他们能对那无所需要的天主有所给予，在我看来，与那些刻意将同样尊荣加给没有感觉、因而无法享受这些尊荣之物的，毫无区别。\par

% structure: p0011 source=h2 target=section reason=relative-heading-rank; base=h2

\section{犹太人的其他规条}

但是，关于他们对食物的过分讲究，对安息日的迷信，对割损礼的夸耀，对斋戒和新月的幻想——这些都完全可笑，不值一提——我认为你不需要从我这里学习什么。因为，接受天主为人的使用所造的一些事物为合宜的，而拒绝另一些为无用和多余的——这怎能是合法的呢？而且，妄说天主，好像祂禁止我们在安息日行善——这岂非不敬？并以肉身的割损礼为蒙拣选的证据而自夸，仿佛因此他们特别为天主所爱——这岂非可笑？至于他们观察月与日，\BibleRef{迦 4:10} 仿佛在侍奉星宿和月亮，并按照自己的倾向分配天主的定命和季节的变迁，一些为喜庆，另一些为哀悼——谁会认为这是敬拜的一部分，而不是愚昧的表现呢？所以，我想你已经充分相信，基督徒确实避免了（犹太人和外邦人）共有的虚妄和错误，以及犹太人的多管闲事和虚夸；但你不应指望从任何凡人那里学到他们独特敬拜天主的奥秘。\par

% structure: p0013 source=h2 target=section reason=relative-heading-rank; base=h2

\section{第五章　基督徒的生活方式}

因为基督徒既不以国家、语言，也不以所遵守的习俗与其他人区别开来。他们既不居住自己的城邑，也不使用特殊的言语，也不过着任何奇特标记的生活。他们所遵循的行为方式并非由任何好钻研之人的思虑或讨论所设计；他们也不像有些人那样，只宣扬任何纯粹人为的学说。但是，他们居住在各个希腊人的城邑和异邦人的城邑，按照每人所分配的际遇，并遵循当地人在衣、食以及其余日常行为上的习俗，却向我们展示出他们那奇妙且令人惊叹的生活方式。他们居住在自己的家乡，却如同寄居者。作为公民，他们与众人共享一切，却又忍受一切如同外人。每一处异乡对他们而言都是故乡，每一处出生之地都是客地。他们结婚，如同众人一样；他们生育子女，却不弃绝后代。他们共用桌子，却不共用床铺。他们处于肉身之内，却不随从肉身而活。\BibleRef{格后 10:3} 他们度日在地上，却是天上的公民。\BibleRef{斐 3:20} 他们遵守制定的法律，同时以其生活超越法律。他们爱所有人，却被所有人迫害。他们不为人知又被定罪；他们被处死，却又重获生命。\BibleRef{格后 6:9} 他们贫穷，却使许多人富足；\BibleRef{格后 6:10} 他们一无所有，却又样样丰盛；他们受羞辱，却在羞辱中得荣耀。他们被诽谤，仍得称义；他们被辱骂，却祝福；\BibleRef{格后 4:12} 他们被侮辱，却以尊敬回报侮辱；他们行善，却作为恶人受罚。受罚时，他们如同被激发生命而喜乐；他们被犹太人视为外人而攻击，被希腊人迫害；然而恨他们的人却无法为自己的仇恨提出任何理由。\par

% structure: p0015 source=h2 target=section reason=relative-heading-rank; base=h2

\section{第六章　基督徒与世界的关系}

简言之，灵魂在身体内怎样，基督徒在世界中也怎样。灵魂分散于身体的所有肢体，基督徒分散于世界所有的城邑。灵魂居于身体之内，却不属于身体；基督徒居于世界之内，却不属于世界。不可见的灵魂由可见的身体护卫，基督徒固然被知道存在于世界，但他们的敬虔却是不可见的。肉身恨恶灵魂，与之交战，\BibleRef{伯前 2:11} 虽然本身未受伤害，只因被阻止享受快乐；世界也恨恶基督徒，虽未受丝毫伤害，只因他们弃绝快乐。灵魂爱那恨恶它的肉身，也爱众肢体；基督徒同样爱那些恨他们的人。灵魂被囚禁在身体内，却保全那身体；基督徒被限制在世界中如同在监牢里，却保全了世界。不朽的灵魂居住在必朽的帐幕中；基督徒如同旅客居住在可朽坏的身体内，期盼天上不朽的居所。灵魂在饮食不足时反而变得更好；同样，基督徒虽每日受罚，人数却越发增多。天主派定了他们这光辉的地位，他们弃之是不合法的。\par

% structure: p0017 source=h2 target=section reason=relative-heading-rank; base=h2

\section{第七章　基督的显现}

因为，正如我所说，交付给他们的并非凡俗的发明，也不是纯粹人类意见的体系，他们判断应当如此谨慎地保存，也没有纯粹的属人奥秘的职务委托给他们，而是确实天主自己，全能者，万物的创造者，不可见的，从天上派遣了并放在人中［那位］真理，圣洁而不可理解的圣言，并坚固地建立在他们的心中。祂并没有，如同人可以想象的，派遣给人们任何仆人、天使、统治者，或任何执掌属地事物者，或任何受委托管理天上事物者，而是万物的创造者和塑造者——祂借着祂创造了诸天——祂将海洋限制在适当的界限内——众星忠实地遵守祂的律令——太阳从祂接受了每日行程的度量——月亮服从祂，受命在夜晚照耀，星星也服从祂，跟随月亮运行；万物都由祂安排，放置在适当的界限内，万物都服从祂——诸天及其中的，大地及其中的，海洋及其中的——火、空气和深渊——高处、深处及中间的一切。祂派遣了这位［使者］给他们。那么，如同可以设想的，是为了施行暴政，或引起恐惧和战栗吗？绝非如此，而是出于仁慈与温和。正如国王差遣他的儿子，一位也是国王，祂也这样差遣了祂；作为天主差遣了祂；作为对人差遣了祂；作为救主差遣了祂，并寻求说服而非强迫我们；因为强制在天主的性格中没有位置。作为召唤我们派遣祂，而非报复地追逐我们；作为爱我们派遣祂，而非审判我们。因为祂还将派遣祂来审判我们，谁能承受祂的出现？……你们没有看见他们被投入野兽之中，以说服他们否认主，却仍不被征服吗？你们没有看见越多的他们被惩罚，其余的人就变得越多吗？这似乎不是人的工作：这是天主的德能；这些是祂显现的证据。\par

% structure: p0019 source=h2 target=section reason=relative-heading-rank; base=h2

\section{第八章 在圣言来临之前人类的悲惨境况}

因为，在祂来临之前，有人理解天主是什么吗？你是否接受那些被认为是值得信赖的哲学家的虚妄愚蠢的学说？其中有人说火是天主，称那火为天主，而他们自己不久也将归于火；有的说水；有的说其他由天主形成的元素。但如果这些理论中有任何一个值得赞许，那么其余被造物中的每一个也可能被宣称为天主。但这样的声明不过是欺骗者令人震惊和错误的言论；没有人见过祂，或使人认识祂，但祂启示了自己。祂借着信德彰显了自己，唯独信德被赐予看得见天主。因为天主，万物的主和塑造者，祂创造了万物，并给它们分派了各自的位置，证明了自己不仅是人类的朋友，而且对他们的表现是长期忍耐的。是的，祂一直是这样的性情，现在仍是，将来永远是，仁慈、善良、不发怒、真实，唯一［绝对］善良的那位；\BibleRef{玛 19:17}祂在心里形成了一个伟大而不可言喻的概念，只与祂的圣子沟通。那么，当祂保持并保守自己的智慧计划隐藏时，祂似乎忽略了我们对祂没有关心。但在祂借着祂的爱子揭示并打开了从起初就已经预备好的事物之后，祂立刻将一切祝福都赐给了我们，使我们既分享祂的恩惠，又能看见并活跃于祂的服务中。我们中有谁曾预料到这些事？祂与祂的圣子一起，按照他们之间的关系，在自己的心中知晓了一切。\par

% structure: p0021 source=h2 target=section reason=relative-heading-rank; base=h2

\section{第九章 为什么圣子如此晚才被派遣}

所以，只要先前的时期延续，祂就容许我们被不受约束的冲动所拖曳，被欲望的快乐和各种情欲所牵引。这并非因为祂喜悦我们的罪，而是祂单纯地容忍它们；也不是因为祂赞同那时行不义的时期，而是祂要塑造一个认识正义的心灵，以便我们在那时期被说服，凭自己的行为不配获得生命，如今却藉着天主的仁慈，得以赐予我们；并且祂显明了，凭我们自己无法进入天主的国，我们却可以藉着天主的能力变得能够。但当我们的邪恶达到顶点，并且清楚显明它的报酬——惩罚和死亡——正临到我们时；当时间已到，天主先前所定要显示祂自己的仁慈和能力的时候，那唯一的天主之爱，因着对人的无限关怀，并没有以仇恨对待我们，也没有推开我们，更没有记念我们的不义来反对我们，反而显出极大的宽忍，担待了我们；祂亲自承担了我们不义的负担，将祂自己的圣子作为赎价给了我们——圣洁者给了悖逆者，无瑕者给了邪恶者，正义者给了不义者，不朽者给了可朽者，永生者给了必死者。因为除了祂的义德，还有什么能遮盖我们的罪呢？除了天主的独生子，我们这些邪恶和不敬神的人，还能藉着谁成义呢？哦，甜蜜的交换！哦，不可测度的作为！哦，超越一切期望的恩惠！许多人的邪恶竟能隐藏在一位正义者内，而这一位的正义竟能使许多悖逆者成义！因此，祂既在先前的时期使我们确信，我们的本性无法获得生命，如今又揭示了那能拯救甚至那些〔先前〕不可能被拯救者的救主，藉着这两件事，祂要引导我们信赖祂的仁慈，视祂为我们的养育者、圣父、导师、谋士、医治者、我们的智慧、光、荣耀、光荣、德能和生命，使我们不再为衣服和食物忧虑。\par

% structure: p0023 source=h2 target=section reason=relative-heading-rank; base=h2

\section{第十章  信仰将带来的福乐}

如果你也渴望拥有这信德，你同样将首先接受对圣父的认识。因为天主爱人，祂为人创造了世界，将世上万物都置于人的权下，赐给人理智和领悟，唯独使人能举目仰望祂，按自己的肖像造了人，差遣祂的独生子到人这里来，应许给那些爱祂的人天上的国，并要赐予他们。当你获得这认识时，你想你会充满何等的喜乐呢？或者，你将怎样爱那先如此爱了你的祂呢？如果你爱祂，你将成为祂仁慈的效法者。不要惊奇，一个人能成为天主的效法者。如果他愿意，就能。因为幸福不在于统治邻人，或寻求压制那些较弱者，或富有而对那些低下者施暴；没有人能藉着这些成为天主的效法者。但这些事丝毫不构成祂的尊威。相反，那承担邻人负担的；那在任何方面优越而乐意帮助那有欠缺者的；那将从天主领受的一切，分施给有需要的人，从而成为那些领受者的神——这样的人是天主的效法者。那时，你虽然还在地上，必将看见天主在天上统治万物；那时，你将开始讲述天主的奥秘；那时，你将爱慕那些因不肯否认天主而忍受惩罚的人；那时，你将谴责世界的欺骗和错误，当你知道什么是真正在天上生活时，当你轻视那在此被视为死亡的事时，当你惧怕那真正的死亡，即为那些将被判入永火的人所存留的、将折磨那些被交付其中的人直到终结的火时。那时，你将钦佩那些为义德的缘故忍受那一时之火的人，并视他们有福，当你认识那火的性质时。\par

% structure: p0025 source=h2 target=section reason=relative-heading-rank; base=h2

\section{第十一章  这些事值得认识和相信}

我不是在谈论我所不熟悉的事，也不追求任何与正当理性相悖的事；但作为宗徒的门徒，我成了外邦人的教师。我将交付给我的事物，传递给那些堪当真理的门徒。因为那由慈爱的圣言正确教导并重生的人，谁不会渴望准确学习圣言向祂的门徒清楚显明的事呢？圣言显现并启示他们，明确地［对他们］说话，不为不信者所理解，却与门徒交谈；他们被祂视为忠实，从而获得了圣父奥秘的知识。为此，祂派遣了圣言，使祂向世界显现；祂虽被［犹太］百姓轻蔑，但当宗徒们宣讲时，却被外邦人所信。这就是那从起初就有的，祂似乎新出现，却被发现是古老的，且常在圣人们的心中重生。这就是那从永远而有的，今日被称为圣子；藉着祂，教会丰富，恩宠广施，在圣人们中增长，提供理解，启示奥秘，宣告时期，为信者喜乐，赐予寻求者；藉着祂，信德的界限不被冲破，先父们设立的界限不被跨越。于是，对法律的敬畏被歌颂，先知们的恩宠被认识，福音的信德得以确立，宗徒的传统得以保存，教会的恩宠欢然雀跃；你若不让这恩宠忧伤，你必能知道圣言所教导的事，按祂的旨意和祂喜悦的时候。凡因那命令我们的圣言的旨意而感动我们所说的一切，我们都以辛劳和对所启示之事的爱传达给你们。\par

% structure: p0027 source=h2 target=section reason=relative-heading-rank; base=h2

\section{第十二章 知识对真正属灵生活的重要性}

当你阅读并仔细聆听这些事后，你必将知道天主赐予那些正确爱祂的人什么，（你们自己）成为喜乐的乐园，在你们内呈现一棵结满各样果实、繁荣兴旺、装饰着各种果子的树。因为在这地方，知识树和生命树都被栽种了；但毁坏人的不是知识树，而是悖逆造成毁灭。那些经文也绝非毫无意义，记载了天主起初将生命树栽种在乐园当中，藉着知识启示通往生命的道路，而当初受造的人没有正确运用这［知识］，就因蛇的欺诈而赤裸。因为没有知识，生命无法存在；没有生命，知识也不稳固。因此两者被栽种在彼此近旁。宗徒察觉了［这种结合］的力，并责备那种没有真道却影响生活的知识，宣称：\BibleQuote{知识使人自高自大，但爱德却能建造人。}因为那以为知道些什么却缺乏真知识、且未经生命所见证的人，其实一无所知，反被蛇欺骗，如同不爱生命一样。但那将知识与敬畏结合，并寻求生命的人，在希望中栽种，等待果实。让你的心作你的智慧，让你的生命成为内心所接受的真正知识。肩负这棵树，展示其果实，你将常收获那些天主所渴望的事物，是蛇不能触及、欺骗无法靠近的；那时\Person{厄娃}不再被败坏，反被信赖如同贞女；救恩得以彰显，宗徒们充满智慧，主的逾越节临近，歌咏团体聚集，并排列有序，圣言因教导圣人而喜乐——藉着圣言，圣父受光荣：愿光荣归于祂，直到永远。亚孟。\par

\SourceRecord{Translated by Alexander Roberts and James Donaldson.；Ante-Nicene Fathers；Vol. 1.；Edited by Alexander Roberts, James Donaldson, and A. Cleveland Coxe.；Buffalo, NY: Christian Literature Publishing Co.,；1885.；<http://www.newadvent.org/fathers/0101.htm>.}
